% BEGIN VOL08-EXPANSION VIII18
\section{Supplementary worked examples}
\label{sec:viii18-pedagogy}

\begin{example}[Cellular homology of the sphere]\label{ex:viii18-ped-01}
With one 0-cell and one $n$-cell, the cellular chain complex of $S^n$ has $C_n=C_0=\mathbb Z$ and no adjacent nonzero groups, so all differentials vanish and $H_n\cong H_0\cong\mathbb Z$.
\end{example}

\begin{example}[Cellular homology of the torus]\label{ex:viii18-ped-02}
For the one-vertex, two-edge, one-face CW structure of $T^2$, both cellular differentials vanish: $d_1=0$, and the commutator attaching word has zero total degree on each 1-cell. Thus $H_2\cong\mathbb Z$, $H_1\cong\mathbb Z^2$, $H_0\cong\mathbb Z$.
\end{example}

\begin{example}[Projective-space differential]\label{ex:viii18-ped-03}
In the one-cell-per-dimension CW structure of $\mathbb{RP}^n$, cellular differentials alternate between multiplication by $2$ and $0$ (with convention depending on dimension indexing), producing the characteristic $\mathbb Z/2$ torsion.
\end{example}

\section{Graded supplementary exercises}
The following exercises progress from direct computation and construction to proof, hypothesis testing, application, and synthesis.

\subsection*{Standard computations and constructions}

\begin{exercise}[Cellular chain rank]\label{exr:viii18-ped-01}
What is the rank of $C_n^{\mathrm{cell}}(X)$ for a CW complex with $c_n$ $n$-cells?
\end{exercise}
\begin{hint}
One generator per cell.
\end{hint}
\begin{solution}
$C_n^{\mathrm{cell}}(X)\cong\mathbb Z^{c_n}$ for integral coefficients.
\end{solution}

\begin{exercise}[Sphere differential]\label{exr:viii18-ped-02}
Why are cellular differentials zero in the two-cell CW structure of $S^n$?
\end{exercise}
\begin{hint}
There are no cells in dimensions $n-1$ or $1$ when $n>1$.
\end{hint}
\begin{solution}
The target/source adjacent chain groups vanish, so the differentials are zero (with the circle case also having $d_1=0$ because the single edge has the same initial and terminal vertex).
\end{solution}

\begin{exercise}[Torus ranks]\label{exr:viii18-ped-03}
State the cellular chain ranks for the standard torus CW structure.
\end{exercise}
\begin{hint}
Count cells $1,2,1$.
\end{hint}
\begin{solution}
$\operatorname{rank}C_0=1$, $\operatorname{rank}C_1=2$, $\operatorname{rank}C_2=1$.
\end{solution}

\begin{exercise}[Degree boundary]\label{exr:viii18-ped-04}
If one 2-cell attaches to one 1-cell by degree $m$, what is $d_2$?
\end{exercise}
\begin{hint}
Cellular incidence is the attaching degree.
\end{hint}
\begin{solution}
$d_2:\mathbb Z\to\mathbb Z$ is multiplication by $m$.
\end{solution}

\begin{exercise}[Euler from cells]\label{exr:viii18-ped-05}
For a finite CW complex, express Euler characteristic using cellular chain ranks.
\end{exercise}
\begin{hint}
Alternate the cell counts.
\end{hint}
\begin{solution}
$\chi(X)=\sum_n(-1)^n\operatorname{rank}C_n^{\mathrm{cell}}=\sum_n(-1)^nc_n$.
\end{solution}

\subsection*{Proofs}

\begin{exercise}[Cellular chains from relative homology]\label{exr:viii18-ped-06}
Explain why $C_n^{\mathrm{cell}}(X)\cong H_n(X^n,X^{n-1})$ is free on the $n$-cells.
\end{exercise}
\begin{hint}
Collapse $X^{n-1}$ and obtain a wedge of $n$-spheres.
\end{hint}
\begin{solution}
$X^n/X^{n-1}\cong\bigvee S^n$ over the $n$-cells; reduced homology in degree $n$ is one copy of $\mathbb Z$ per sphere/cell.
\end{solution}

\begin{exercise}[$d^2=0$]\label{exr:viii18-ped-07}
Why do cellular boundary maps satisfy $d_{n-1}d_n=0$?
\end{exercise}
\begin{hint}
They arise as connecting maps in exact sequences of skeletal triples/pairs.
\end{hint}
\begin{solution}
Naturality and exactness of the long exact sequences of the filtration make consecutive connecting homomorphisms compose to zero, producing a chain complex.
\end{solution}

\begin{exercise}[Cellular equals singular]\label{exr:viii18-ped-08}
State why cellular homology computes ordinary singular homology of a CW complex.
\end{exercise}
\begin{hint}
Use the skeletal filtration spectral/exact-sequence argument.
\end{hint}
\begin{solution}
The filtration yields a cellular chain complex whose homology is naturally isomorphic to $H_*(X)$; only relative groups in the matching cell dimension survive.
\end{solution}

\begin{exercise}[Euler-Poincare]\label{exr:viii18-ped-09}
Prove the alternating sum of cell counts equals the alternating sum of homology ranks for a finite CW complex.
\end{exercise}
\begin{hint}
Use a finite free chain complex and cancellation of boundary ranks.
\end{hint}
\begin{solution}
Writing $\operatorname{rank}C_n=\operatorname{rank}B_n+\operatorname{rank}H_n+\operatorname{rank}B_{n-1}$, the boundary-rank terms cancel in the alternating sum.
\end{solution}

\subsection*{Counterexamples and hypothesis tests}

\begin{exercise}[Need triangulation]\label{exr:viii18-ped-10}
Test: cellular homology requires subdividing every CW cell into simplices.
\end{exercise}
\begin{hint}
It works directly from CW skeleta.
\end{hint}
\begin{solution}
False. The point is to compute from cells without triangulating them.
\end{solution}

\begin{exercise}[Boundary is cell count only]\label{exr:viii18-ped-11}
Test: cellular boundary matrices are determined solely by the number of cells.
\end{exercise}
\begin{hint}
Attaching maps determine incidence degrees.
\end{hint}
\begin{solution}
False. Equal cell counts can have different attaching maps and different homology.
\end{solution}

\begin{exercise}[All CW differentials zero]\label{exr:viii18-ped-12}
Test: cellular differentials vanish because cells are open disks.
\end{exercise}
\begin{hint}
The attaching maps carry the topology.
\end{hint}
\begin{solution}
False. Nonzero attaching degrees give nonzero differentials and torsion.
\end{solution}

\subsection*{Applications and investigations}

\begin{exercise}[Sparse topology]\label{exr:viii18-ped-13}
Why is cellular homology often much smaller than simplicial homology for the same space?
\end{exercise}
\begin{hint}
A CW model can use very few cells.
\end{hint}
\begin{solution}
Each cell contributes one chain generator, so a minimal CW decomposition can replace a huge triangulation by tiny boundary matrices.
\end{solution}

\begin{exercise}[Attachment diagnostics]\label{exr:viii18-ped-14}
How can a cellular boundary matrix help diagnose whether a proposed CW attachment encoded the intended relation?
\end{exercise}
\begin{hint}
Its entries are signed degrees/incidences of attaching maps.
\end{hint}
\begin{solution}
Unexpected matrix entries reveal incorrect winding/orientation of attaching maps; homology then provides a global consistency check.
\end{solution}

\subsection*{Challenge problems}

\begin{exercise}[Projective torsion]\label{exr:viii18-ped-15}
Explain qualitatively how a multiplication-by-2 cellular differential creates $\mathbb Z/2$ homology in projective space.
\end{exercise}
\begin{hint}
Take a cokernel of $2:\mathbb Z\to\mathbb Z$ when the adjacent differential is zero.
\end{hint}
\begin{solution}
At the relevant degree, cycles are all of $\mathbb Z$ while boundaries are $2\mathbb Z$, so the quotient is $\mathbb Z/2$.
\end{solution}

\begin{exercise}[Cellular synthesis]\label{exr:viii18-ped-16}
Describe the complete workflow from a finite CW description to homology.
\end{exercise}
\begin{hint}
Count cells, determine attaching degrees, build matrices, compute kernels/images.
\end{hint}
\begin{solution}
Choose orientations for cells; form one free generator per cell; compute cellular differentials from attaching-map degrees; verify $d^2=0$; reduce the integer matrices (e.g. Smith form) to obtain kernels modulo images and hence all homology groups.
\end{solution}

% END VOL08-EXPANSION VIII18
